\section{Selected Composition and Integrated Verification}
\label{sec:verification}
\subsection{What Was Selected, and What Remains Experimental}
The local final manifest selects BM25 retrieval, ambiguity-aware evidence
selection, deterministic factual assembly, guarded model planning and persistent
in-app outreach. The exact profile is listed in Appendix~\ref{sec:software}. Table~\ref{tab:design-decisions} distinguishes
comparison-based choices from requirements that the implementation must enforce.

The experiments do not identify a single composition that wins on every
criterion. The evidence-interface comparison tests a development evidence interface; the fresh factual comparison supports
the retained factual fallback. The action-planning and model-allocation comparisons compare planners using supplied
state cards, whereas the product adapter still supplies delivery/event proxies.
The learner/timing simulation evaluates estimator/timing alternatives outside the tutoring service.
The live-model integration trial exercises an integrated development configuration with actual model
calls, and the goal-completion regression isolates a later lifecycle correction without model calls.
Thus neither the component improvements nor the historical integrated run
constitute an end-to-end qualification of all current components together.


\begin{table}[htbp]
\centering\small
\caption{Current disposition of the principal designs.}
\label{tab:design-decisions}
\begin{tabularx}{\linewidth}{@{}L{0.23\linewidth}Y Y@{}}
\toprule Design boundary & Current disposition & Basis and unresolved issue \\\midrule
Source/policy binding & Retain reviewed, immutable releases & Required for authority and reproducible dialogue; supported by publication tests, not a performance ranking. \\
Retrieval and factual wording & Retain BM25/dominance and deterministic assembly & The \hyperref[study-B]{fresh factual comparison} supports the simpler local fallback, not a qualified winner; grounded and boundary scores remain below targets. \\
Question-to-evidence mapping & Retain as development design & The \hyperref[study-A]{evidence-interface comparison} improves coverage; not all gates pass and not a release-wide 89.42\% claim. \\
Autonomous planning/wording & Retain guarded planning with Luna for both planning and wording & Synthetic utility supports progression to integration, despite lower preferred-action agreement; product proxy inputs and individualized instruction remain unvalidated. \\
Learner estimator/timing & Keep current evidence counters; BKT/value remains experimental & The \hyperref[study-E]{learner/timing simulation} supports further work, but its estimator is not wired into the default planning adapter. \\
Goal lifecycle correction & Keep objective-scoped completion in the working tree & The \hyperref[study-H]{goal-completion regression} removes unsupported completions; full frozen-release requalification is separate. \\
Flexible generation and vision & Keep unsuccessful candidates as development evidence & Critical semantic/attribution failures or no end-to-end gain prevented promotion. \\
Delivery and persistence & Retain leased jobs, authority rechecks and stable delivery keys & Required for consent and restart behaviour; local operational checks bound the evidence. \\
\bottomrule
\end{tabularx}
\end{table}

The selected profile uses GPT-5.6 Luna for bounded complex planning and wording
strategies. Flexible instructional trials are separate configurations. Likewise,
the current objective-scoped correction postdates the frozen local qualification:
its experimental component profile and matched regression study support that
correction, not a claim that the complete final release was requalified.

\subsection{Thirty Virtual Days with Actual Model Calls}
\label{sec:provider-study}
The \hyperref[study-G]{live-model integration trial} executes the integrated development composition over 24 synthetic
histories and 30 virtual days. It uses actual external LLM calls for its
configured model pathways: 654 calls were recorded, 643 completed and 11
reached the output limit. The run processed 484 tutor turns, delivered 16 proactive
messages and exercised 24 service restarts. Consent was disabled on virtual
days 10--19; no proactive messages were delivered in that interval.

The software actually executed dialogue processing, persistent learner records,
autonomous jobs, delivery and restart recovery. Student utterances, attendance and receptivity were synthetic, and time
advanced virtually. This operational trial did not score changes in mastery
and does not establish 30 days of uptime. The separate learner/timing simulation
and goal-completion regression make no external model calls; they examine
simulated learning mechanisms and state-management correctness, respectively.

Thirteen tutoring turns ended in a safe graph failure, including two invalid
question-focus selections. A post-run AI diagnostic reviewed 48 endpoint
responses, all 13 failures and all 16 check-ins: 77 deliberately selected cases,
not a representative estimate of response quality. All 16 check-ins used a
generic wrapper around an entire source card; their effect on students was
not measured. In the saved Socratic first-turn
case, the tutor asks for the student's current explanation but does not make
the question specific to the concept's underlying difficulty
(Appendix~\ref{sec:profile-examples}).

The operational finding is that the composed software can initiate and preserve
support under the tested controls. The instructional finding is narrower:
successful authorization and delivery did not ensure a useful, personalized
question. Provider cost, latency and the precise historical configuration remain
in the linked study record (Appendix~\ref{sec:study-index}).

\subsection{A Goal-Completion Defect and Its Revision Cycle}
\label{sec:goal-study}
A service reproduction exposed a missing invariant: completing goal $g$ must
depend on evidence for $C(g)$, not the strongest concept anywhere in the learner
record. With cache-coherence and virtual-memory goals active, two correct
cache-coherence attempts incorrectly completed both goals; reopening SQLite
preserved the error. The old service shortcut and goal interpreter both used
evidence too broadly.

The revised design applies Equation~\ref{eq:goal-completion} separately to each
goal and accepts only committed state from the relevant scope. Multi-goal and
multi-concept regressions cover missing, unrelated, incorrect, ambiguous and
uncommitted evidence, plus restart and continued work on another goal. The
prospective \hyperref[study-H]{goal-completion regression} then compares the old and corrected implementations.
Only two runtime source files differ between the archived arms; the driver,
six concepts, seeds, profiles and dependencies match.

Each arm runs 72 histories: six personas, two simulator families, three seeds
and reactive/autonomous conditions, with a restart on virtual day 15. The
product service, worker and SQLite paths execute; wording and planning are
deterministic. Shared initial conditions and seeds do not imply identical
later transcripts, because changed interventions can change simulated replies.

\begin{table}[htbp]
\centering\small
\caption{The \hyperref[study-H]{goal-completion regression}: matched lifecycle correction, 72 histories per arm.
The last two rows concern the 36 autonomous histories per arm.}
\label{tab:goal-correction}
\begin{tabularx}{\linewidth}{@{}Yrr@{}}
\toprule Measure & Original & Objective-scoped \\\midrule
Completed goals & 26 & 15 \\
Supported by target concepts & 14 & 15 \\
Unsupported completion & 12 & 0 \\
Histories with unsupported completion & 11 & 0 \\
Consistent restarts & 72/72 & 72/72 \\
Attribution / assessment agreement & 100\% / 100\% & 100\% / 100\% \\
Quiet-hour / frequency / cooldown violations & 0 / 0 / 0 & 0 / 0 / 0 \\\midrule
Mean proactive messages & 9.056 & 10.611 \\
Mean final simulated mastery & 0.3632 & 0.3608 \\
\bottomrule
\end{tabularx}
\end{table}

The old assessment gates passed even when goal completion was wrong. The
corrected arm passed all 19 run gates. Target-supported completions numbered
14 in the original arm and 15 in the corrected arm; unsupported completions
numbered 12 and zero. Total completed goals measure software status decisions,
not student learning success. A separate audit matches each completion to committed
target-concept evidence. Equal virtual timestamps cannot establish exact commit
ordering, so service regressions separately test that boundary.

The correction increased proactive delivery by 56 messages across the autonomous
histories. Mean simulated mastery changed by $-0.0024$, and mean wasted-message
fraction rose from $0.3496$ to $0.3572$. The evidence supports \emph{Keep} for
objective-scoped state correctness. The mastery difference is a descriptive
secondary outcome relative to the previous implementation, not a measurement
of learning effects in students. Eleven autonomous histories increased, nine
decreased and sixteen were unchanged. This regression uses previously exposed
concepts.

A retrospective audit of the earlier 72-history multi-concept trial
found ten unsupported completions among 22 completed goals in eight histories.
Its original measurements and prior successful assessment correction are
retained, but they cannot establish correctness of the now-revised lifecycle.
The four completed goals in the \hyperref[study-G]{live-model integration trial} showed no such mismatch; that small audit
does not qualify all provider-backed paths. Appendix~\ref{sec:historical-autonomy}
preserves the original simulated outcome and its later interpretation.

\subsection{Verification and Acceptance Boundaries}
This revision cycle supplies a concrete software development lifecycle (SDLC) argument: reproduce the defect,
state the missing invariant, change the responsible boundary, compare the
implementations, and retain both the failed baseline and bounded acceptance.
The focused service/component/audit suite passed 72 tests, including 30 new
regressions. The broader Python run and its targeted rechecks are documented in
Appendix~\ref{sec:sdlc}; they are not a single clean post-correction full run.

For course configuration and human control, browser diagnostics exercised preview, approval, reload and
withdrawal. Local qualification recorded 43 application checks and 51 container
regressions; restart, restore and rollback were tested under that environment.
These support publication and operational contracts. They do not replace factual acceptance, pedagogical review or real-course
validation of continuing support.
At submission, course configuration and human control have bounded local
verification; factual support remains below acceptance targets; teaching-style
alignment lacks validation with the actual instructor; and continuing support
has operational and lifecycle evidence but not demonstrated learning
benefit. The prototype implements the course and tutoring workflows, but has not met
all five acceptance requirements: factual-quality gates remain unmet, and
instructor fidelity and real-student learning effects remain unvalidated.
